% Awesome Source CV LaTeX Template
%
% This template has been downloaded from:
% https://github.com/darwiin/awesome-neue-latex-cv
%
% Author:
% Christophe Roger
%
% Template license:
% CC BY-SA 4.0 (https://creativecommons.org/licenses/by-sa/4.0/)

%Section: Project
\sectionTitle{Projects}{\faLaptop}
\begin{projects}
			 % \project {YAAC Another Awesome CV}{2013 - 2018} {\github{darwiin/yaac-another-awesome-cv} \website{https://www.overleaf.com/latex/templates/awesome-source-cv/wrdjtkkytqcw}{Template sur Overleaf}} {Template \LaTeX pour Curiculum Vitæ utilisant les icônes \href{https://fontawesome.com}{Font Awesome} et la police de caractère \href{https://fonts.google.com/specimen/Source+Sans+Pro}{Adobe Source Sans Pro}. YAAC Another Awesome CV a d'abord été créé comme un template simple pour CV à vocation technologique.} {\LaTeX,Sublime Text}
 
    \project{uniMedic}{2020 - 2021}{\github{uniMedic}}{A medical chatbot was developed using Deep Learning for real-time diagnostics. The mobile application was created with React Native, the backend with Express.js to manage records and medical histories, and Flask to consume Artificial Intelligence models. Deployed on the AWS cloud using Amplify, Cognito, and DynamoDB.} {React Native, Express.js, Flask, AWS, Pytorch}
    
    \project{RasPi-BloodView}{2020 - 2021}{\github{Enigma-A-I/RasPi-BloodView}} {Blood cell detection system deployed on Raspberry Pi 4 using PyTorch. The convolutional neural network model allows quick and cost-effective recognition of 3 types of blood cells. Features a custom HTML/CSS interface rendered with Flask and a file management system for secure storage on the server.} {PyTorch, Raspberry Pi}
    
    % \project{faceAcecom}{2020}{\github{AcecomFCUNI/faceAcecom}} {The faceAcecom project implements a facial identification and verification system for ACECOM-IA members. It combines machine learning techniques such as HOG and SVM with convolutional neural networks for deep learning in PyTorch and TensorFlow, allowing real-time facial detection. It was implemented in the cloud to provide scalable facial recognition services.} {Machine Learning, Deep Learning, Pytorch, Tensorflow}
\end{projects}